% =====================================================================
%  2bat-ud6-10.tex  —  SESSIÓ 10: Llegir una fitxa tècnica d'enquesta
%  (sense preàmbul: s'inclou des de main.tex)
% =====================================================================
\horatitol{Sessió 10 — Llegir una fitxa tècnica d'enquesta}%
{Aprendre a llegir la fitxa tècnica d'una enquesta: quins elements té i quina informació cal mirar sempre per valorar-ne la fiabilitat.}

\subsection{Idea clau}

Ensenyem a llegir la \textbf{fitxa tècnica} d'una enquesta: el document, sovint amagat al
final dels estudis, que conté la informació clau per valorar-ne la fiabilitat. Aprendre a
buscar-la i interpretar-la és una rutina de lectura crítica essencial.

\begin{idea}
\textbf{Anatomia d'una fitxa tècnica.} Els elements que cal mirar sempre:
\begin{itemize}[leftmargin=1.4em, itemsep=2pt]
  \item \textbf{Població estudiada:} a qui es vol conèixer (el col·lectiu de referència).
  \item \textbf{Mida de la mostra:} quantes persones s'han enquestat.
  \item \textbf{Mètode de selecció:} com s'ha triat la mostra (a l'atzar? per quotes? qui
        podia respondre?).
  \item \textbf{Dates del treball de camp:} quan s'ha fet (un resultat pot quedar desfasat).
  \item \textbf{Marge d'error} i \textbf{nivell de confiança:} la precisió i la fiabilitat
        de l'estimació.
\end{itemize}
\textbf{La rutina de lectura crítica:} abans de fer cas d'una enquesta, comprova qui ha
respost, quants, com s'han triat i quin marge d'error té.
\end{idea}

\noindent\textbf{Exemple de fitxa tècnica.}
\begin{center}
\renewcommand{\arraystretch}{1.3}
\begin{tabular}{l l}
\toprule
\textbf{Element} & \textbf{Valor} \\
\midrule
Població & Persones majors de 18 anys d'un municipi \\
Mida de la mostra & $800$ persones \\
Selecció & Aleatòria, per telèfon (fix i mòbil) \\
Treball de camp & Març de 2027 \\
Marge d'error & $\pm 3,5\%$ \\
Nivell de confiança & $95\%$ \\
\bottomrule
\end{tabular}
\end{center}

\subsection{Exemples}

\begin{exemple}[llegir la fitxa element a element]
De la fitxa anterior en llegim: l'estudi vol conèixer els \textbf{majors de 18 anys} del
municipi (població); ha enquestat \textbf{800} persones (mostra) triades \textbf{a l'atzar}
per telèfon (mètode raonable, arriba a fix i mòbil); el treball de camp és de \textbf{març
de 2027} (cal comprovar que no hagi quedat desfasat); i té un \textbf{marge de $\pm 3,5\%$}
amb \textbf{confiança del $95\%$}. En conjunt, sembla una enquesta \textbf{fiable}.
\end{exemple}

\begin{exemple}[una fitxa que aixeca sospites]
Una altra fitxa diu: mostra de \textbf{80} persones, selecció «\textbf{voluntària per web}»,
sense marge d'error. Aquí hi ha diverses \textbf{alertes}: mostra petita, selecció
\textbf{autoseleccionada} (biaix), i cap marge d'error (no expressa la incertesa). Els seus
resultats s'han de prendre amb molta cautela.
\end{exemple}

\subsection{Exercicis resolts}

\begin{exresolt}[valorar la fiabilitat d'una fitxa]
Una fitxa indica: població «usuaris d'un servei», mostra $500$, selecció aleatòria, marge
$\pm 4\%$, confiança $95\%$. Valora'n la fiabilitat.
\solucio
Sembla \textbf{fiable}: la selecció és \textbf{aleatòria} (evita biaix), la mostra és
raonable ($500$) i expressa la incertesa amb un marge ($\pm 4\%$) i un nivell de confiança
($95\%$). Caldria conèixer les dates del treball de camp per assegurar que no està
desfasada.
\resposta{Fiable; convindria comprovar-ne les dates.}
\end{exresolt}

\begin{exresolt}[detectar la feblesa]
Una fitxa diu: mostra de $120$ persones recollida «\textbf{a la sortida d'un acte del
servei}». Quina feblesa hi detectes?
\solucio
La selecció no és aleatòria: s'enquesta només qui assisteix a un acte del servei (perfil ja
implicat i satisfet). És un \textbf{biaix de selecció} que sobreestimaria la valoració
positiva. A més, la mostra és petita. Els resultats no es poden generalitzar a tots els
usuaris.
\resposta{Biaix de selecció (només assistents) i mostra petita.}
\end{exresolt}

\subsection{Exercicis per fer}

\begin{exercicis}
\begin{exlist}
  \item Enumera els cinc elements que cal mirar sempre en una fitxa tècnica.
  \item Una fitxa indica mostra $1\,000$, selecció aleatòria, marge $\pm 3\%$, confiança
        $95\%$, treball de camp fa tres setmanes. Sembla fiable? Per què?
  \item Una fitxa no indica el marge d'error ni el mètode de selecció. Quina actitud has de
        tenir davant dels seus resultats?
  \item Per què les dates del treball de camp són importants? Posa un exemple de resultat
        que pugui quedar desfasat.
  \item \textbf{(Comparació)} Dues enquestes mesuren el mateix: una amb $n=1\,200$
        aleatòria i marge $\pm 3\%$; l'altra amb $n=150$ voluntària i sense marge. De
        quina et fies més i per què?
  \item \textbf{(Llista de comprovacions)} Escriu la teva pròpia llista de comprovacions
        (les preguntes que faràs sempre) abans de fer cas d'una enquesta.
\end{exlist}
\end{exercicis}

% ---------------------------------------------------------------------
\fullsolucions{Sessió 10}
\begin{solucions}
\begin{sollist}
  \item Població estudiada; mida de la mostra; mètode de selecció; dates del treball de
        camp; marge d'error i nivell de confiança.
  \item Sí, sembla fiable: mostra gran ($1\,000$), selecció aleatòria, marge petit
        ($\pm 3\%$), confiança $95\%$ i dades recents. Compleix tots els criteris clau.
  \item Molta cautela: sense marge d'error no expressa la incertesa i sense el mètode de
        selecció no es pot saber si la mostra és representativa; els resultats no són
        fiables tal com es presenten.
  \item Resposta oberta. Idea: una enquesta reflecteix la situació del moment en què es va
        fer; per exemple, una enquesta d'opinió sobre un servei feta abans d'un canvi
        important pot no reflectir l'opinió actual.
  \item De la de $n=1\,200$ aleatòria: mostra gran, selecció a l'atzar i marge conegut. La
        de $150$ voluntària té biaix d'autoselecció i no expressa la incertesa.
  \item Resposta oberta. Ha d'incloure preguntes com: qui és la població?, quants han
        respost?, com s'han triat (a l'atzar?)?, quin marge d'error i confiança té?, quan
        es va fer?
\end{sollist}
\end{solucions}
